\documentclass[10pt,a4paper]{amsart}
\usepackage{amsmath,amssymb,amsthm,graphicx}
% Retained minimal preamble for the Stacks Project chapter varieties.tex.
% Verbatim excerpt of the chapter preamble: Xy-pic setup (native preamble lines 14-18), the native theorem
% declarations with the mathematical macros (native preamble lines 154-217), and the project own no-op
% definitions of the reference/slogan/history display environments (native wrapper.tex lines 1-3) that the
% counted statement uses for its slogan. Copied mechanically; nothing added, removed or redefined.
% For commutative diagrams we use Xy-pic
\usepackage[all]{xy}

% We use 2cell for 2-commutative diagrams.
\xyoption{2cell}
% Theorem environments.
%
\theoremstyle{plain}
\newtheorem{theorem}[subsection]{Theorem}
\newtheorem{proposition}[subsection]{Proposition}
\newtheorem{lemma}[subsection]{Lemma}

\theoremstyle{definition}
\newtheorem{definition}[subsection]{Definition}
\newtheorem{example}[subsection]{Example}
\newtheorem{exercise}[subsection]{Exercise}
\newtheorem{situation}[subsection]{Situation}

\theoremstyle{remark}
\newtheorem{remark}[subsection]{Remark}
\newtheorem{remarks}[subsection]{Remarks}

\numberwithin{equation}{subsection}

% Macros
%
\def\lim{\mathop{\mathrm{lim}}\nolimits}
\def\colim{\mathop{\mathrm{colim}}\nolimits}
\def\Spec{\mathop{\mathrm{Spec}}}
\def\Hom{\mathop{\mathrm{Hom}}\nolimits}
\def\Ext{\mathop{\mathrm{Ext}}\nolimits}
\def\SheafHom{\mathop{\mathcal{H}\!\mathit{om}}\nolimits}
\def\SheafExt{\mathop{\mathcal{E}\!\mathit{xt}}\nolimits}
\def\Sch{\mathit{Sch}}
\def\Mor{\mathop{\mathrm{Mor}}\nolimits}
\def\Ob{\mathop{\mathrm{Ob}}\nolimits}
\def\Sh{\mathop{\mathit{Sh}}\nolimits}
\def\NL{\mathop{N\!L}\nolimits}
\def\CH{\mathop{\mathrm{CH}}\nolimits}
\def\proetale{{pro\text{-}\acute{e}tale}}
\def\etale{{\acute{e}tale}}
\def\QCoh{\mathit{QCoh}}
\def\Ker{\mathop{\mathrm{Ker}}}
\def\Im{\mathop{\mathrm{Im}}}
\def\Coker{\mathop{\mathrm{Coker}}}
\def\Coim{\mathop{\mathrm{Coim}}}

% Boxtimes
%
\DeclareMathSymbol{\boxtimes}{\mathbin}{AMSa}{"02}

%
% Macros for moduli stacks/spaces
%
\def\QCohstack{\mathcal{QC}\!\mathit{oh}}
\def\Cohstack{\mathcal{C}\!\mathit{oh}}
\def\Spacesstack{\mathcal{S}\!\mathit{paces}}
\def\Quotfunctor{\mathrm{Quot}}
\def\Hilbfunctor{\mathrm{Hilb}}
\def\Curvesstack{\mathcal{C}\!\mathit{urves}}
\def\Polarizedstack{\mathcal{P}\!\mathit{olarized}}
\def\Complexesstack{\mathcal{C}\!\mathit{omplexes}}
% \Pic is the operator that assigns to X its picard group, usage \Pic(X)
% \Picardstack_{X/B} denotes the Picard stack of X over B
% \Picardfunctor_{X/B} denotes the Picard functor of X over B
\def\Pic{\mathop{\mathrm{Pic}}\nolimits}
\def\Picardstack{\mathcal{P}\!\mathit{ic}}
\def\Picardfunctor{\mathrm{Pic}}
\def\Deformationcategory{\mathcal{D}\!\mathit{ef}}
\newenvironment{reference}{}{}
\newenvironment{slogan}{}{}
\newenvironment{history}{}{}
\usepackage{cleveref}
\crefname{originalreference}{original reference}{original references}
\Crefname{originalreference}{Original reference}{Original references}
\makeatletter
\newlabel{morphisms-lemma-finite-proper}{{01WN}{source}{}{sourcepointer-1}{}}
\newlabel{morphisms-lemma-finite-proper@cref}{{[originalreference][0][]01WN}{[1][0][]source}{}{}{}}
\makeatother
\crefname{originallemma}{original lemma}{original lemmas}
\Crefname{originallemma}{Original lemma}{Original lemmas}
\makeatletter
\newlabel{lemma-check-invertible-sheaf-trivial}{{0B40}{source}{}{sourcepointer-2}{}}
\newlabel{lemma-check-invertible-sheaf-trivial@cref}{{[originallemma][0][]0B40}{[1][0][]source}{}{}{}}
\makeatother
\crefname{originallemma}{original lemma}{original lemmas}
\Crefname{originallemma}{Original lemma}{Original lemmas}
\makeatletter
\newlabel{lemma-degree-tensor-product}{{0AYX}{source}{}{sourcepointer-3}{}}
\newlabel{lemma-degree-tensor-product@cref}{{[originallemma][0][]0AYX}{[1][0][]source}{}{}{}}
\makeatother
\crefname{originalreference}{original reference}{original references}
\Crefname{originalreference}{Original reference}{Original references}
\makeatletter
\newlabel{morphisms-lemma-closed-point-fibre-locally-finite-type}{{01TF}{source}{}{sourcepointer-4}{}}
\newlabel{morphisms-lemma-closed-point-fibre-locally-finite-type@cref}{{[originalreference][0][]01TF}{[1][0][]source}{}{}{}}
\makeatother
\crefname{originalreference}{original reference}{original references}
\Crefname{originalreference}{Original reference}{Original references}
\makeatletter
\newlabel{properties-lemma-affine-cap-s-open}{{01PV}{source}{}{sourcepointer-5}{}}
\newlabel{properties-lemma-affine-cap-s-open@cref}{{[originalreference][0][]01PV}{[1][0][]source}{}{}{}}
\makeatother
\crefname{originalreference}{original reference}{original references}
\Crefname{originalreference}{Original reference}{Original references}
\makeatletter
\newlabel{coherent-lemma-vanshing-gives-ample}{{0B5U}{source}{}{sourcepointer-6}{}}
\newlabel{coherent-lemma-vanshing-gives-ample@cref}{{[originalreference][0][]0B5U}{[1][0][]source}{}{}{}}
\makeatother
\crefname{originalreference}{original reference}{original references}
\Crefname{originalreference}{Original reference}{Original references}
\makeatletter
\newlabel{coherent-lemma-coherent-support-dimension-0}{{0B3J}{source}{}{sourcepointer-7}{}}
\newlabel{coherent-lemma-coherent-support-dimension-0@cref}{{[originalreference][0][]0B3J}{[1][0][]source}{}{}{}}
\makeatother
\crefname{originalreference}{original reference}{original references}
\Crefname{originalreference}{Original reference}{Original references}
\makeatletter
\newlabel{coherent-lemma-section-affine-open-kills-classes}{{01XR}{source}{}{sourcepointer-8}{}}
\newlabel{coherent-lemma-section-affine-open-kills-classes@cref}{{[originalreference][0][]01XR}{[1][0][]source}{}{}{}}
\makeatother
\title{An invertible sheaf on a proper curve is ample exactly when its degree is positive}
\author{The Stacks Project Authors}
\date{stacks-project:varieties.tex:local-snapshot-20261011; GNU FDL 1.2 or later}
\begin{document}
\maketitle
\noindent\textit{Editorial note.} All delivered mathematical text is the author's own, copied verbatim from the pinned Stacks Project chapter file varieties.tex; no statement or proof step was rewritten, reordered or removed.\par\medskip
\noindent\textit{Editorial note.} The counted span cites eight cross-chapter labels, each recorded in assembly.external\_references with its canonical tag as a navigation pointer.\par\medskip
\noindent\textit{Editorial note.} The proof cites no bibliography item inside the counted span.\par\medskip
\noindent\textit{Allowed standard prerequisites (editorial).}\begin{itemize}
\item Sign and triviality criteria for invertible sheaves on a proper curve. Stacks Project, Varieties, Lemma lemma-check-invertible-sheaf-trivial (tag 0B40). Turns a nonzero section vanishing at a point into strict positivity of the degree.
\item The Riemann-Roch type formula for the Euler characteristic of powers. Stacks Project, Varieties, Lemma lemma-degree-on-proper-curve (tag 0AYV). Gives chi(X, L\textasciicircum{}n) = n deg(L) + chi(X, O\_X), which drives the growth estimate.
\end{itemize}\medskip
\ifdefined\hypertarget\hypertarget{sourcepointer-1}{}\fi
\noindent\textit{Source reference.} Original reference 01WN: sources/stacks-chow/tags.txt, source lines 2488--2488. Cross-chapter standard prerequisite cited by the author inside the counted span; it is used as a pointer and its own proof is not part of this item.\par\medskip
\ifdefined\hypertarget\hypertarget{sourcepointer-2}{}\fi
\noindent\textit{Source reference.} Original lemma 0B40: sources/stacks-chow/tags.txt, source lines 13983--13983. Cross-chapter standard prerequisite cited by the author inside the counted span; it is used as a pointer and its own proof is not part of this item.\par\medskip
\ifdefined\hypertarget\hypertarget{sourcepointer-3}{}\fi
\noindent\textit{Source reference.} Original lemma 0AYX: sources/stacks-chow/tags.txt, source lines 13805--13805. Cross-chapter standard prerequisite cited by the author inside the counted span; it is used as a pointer and its own proof is not part of this item.\par\medskip
\ifdefined\hypertarget\hypertarget{sourcepointer-4}{}\fi
\noindent\textit{Source reference.} Original reference 01TF: sources/stacks-chow/tags.txt, source lines 2372--2372. Cross-chapter standard prerequisite cited by the author inside the counted span; it is used as a pointer and its own proof is not part of this item.\par\medskip
\ifdefined\hypertarget\hypertarget{sourcepointer-5}{}\fi
\noindent\textit{Source reference.} Original reference 01PV: sources/stacks-chow/tags.txt, source lines 2244--2244. Cross-chapter standard prerequisite cited by the author inside the counted span; it is used as a pointer and its own proof is not part of this item.\par\medskip
\ifdefined\hypertarget\hypertarget{sourcepointer-6}{}\fi
\noindent\textit{Source reference.} Original reference 0B5U: sources/stacks-chow/tags.txt, source lines 14047--14047. Cross-chapter standard prerequisite cited by the author inside the counted span; it is used as a pointer and its own proof is not part of this item.\par\medskip
\ifdefined\hypertarget\hypertarget{sourcepointer-7}{}\fi
\noindent\textit{Source reference.} Original reference 0B3J: sources/stacks-chow/tags.txt, source lines 13967--13967. Cross-chapter standard prerequisite cited by the author inside the counted span; it is used as a pointer and its own proof is not part of this item.\par\medskip
\ifdefined\hypertarget\hypertarget{sourcepointer-8}{}\fi
\noindent\textit{Source reference.} Original reference 01XR: sources/stacks-chow/tags.txt, source lines 2530--2530. Cross-chapter standard prerequisite cited by the author inside the counted span; it is used as a pointer and its own proof is not part of this item.\par\medskip
\begin{lemma}
\label{lemma-ample-curve}
Let $k$ be a field. Let $X$ be a proper curve over $k$.
Let $\mathcal{L}$ be an invertible $\mathcal{O}_X$-module.
Then $\mathcal{L}$ is ample if and only if $\deg(\mathcal{L}) > 0$.
\end{lemma}

\begin{proof}
If $\mathcal{L}$ is ample, then there exists an $n > 0$ and a section
$s \in H^0(X, \mathcal{L}^{\otimes n})$ with $X_s$ affine. Since
$X$ isn't affine (otherwise by
Morphisms, Lemma \ref{morphisms-lemma-finite-proper}
$X$ would be finite), we see that $s$ vanishes at some point.
Hence $\deg(\mathcal{L}^{\otimes n}) > 0$ by
Lemma \ref{lemma-check-invertible-sheaf-trivial}.
By Lemma \ref{lemma-degree-tensor-product}
we conclude that $\deg(\mathcal{L}) = 1/n\deg(\mathcal{L}^{\otimes n}) > 0$.

\medskip\noindent
Assume $\deg(\mathcal{L}) > 0$. Then
$$
\dim_k H^0(X, \mathcal{L}^{\otimes n}) \geq \chi(X, \mathcal{L}^n)
= n\deg(\mathcal{L}) + \chi(X, \mathcal{O}_X)
$$
grows linearly with $n$. Hence for any finite collection of closed
points $x_1, \ldots, x_t$ of $X$, we can find an $n$ such that
$\dim_k H^0(X, \mathcal{L}^{\otimes n}) > \sum \dim_k \kappa(x_i)$.
(Recall that by Hilbert Nullstellensatz, the extension fields
$\kappa(x_i)/k$ are finite, see for example
Morphisms, Lemma \ref{morphisms-lemma-closed-point-fibre-locally-finite-type}).
Hence we can find a nonzero $s \in H^0(X, \mathcal{L}^{\otimes n})$
vanishing in $x_1, \ldots, x_t$. In particular, if we choose
$x_1, \ldots, x_t$ such that $X \setminus \{x_1, \ldots, x_t\}$
is affine, then $X_s$ is affine too
(for example by Properties, Lemma \ref{properties-lemma-affine-cap-s-open}
although if we choose our finite set such that
$\mathcal{L}|_{X \setminus \{x_1, \ldots, x_t\}}$ is trivial, then
it is immediate). The conclusion is that we can find an $n > 0$
and a nonzero section $s \in H^0(X, \mathcal{L}^{\otimes n})$ such
that $X_s$ is affine.

\medskip\noindent
We will show that for every quasi-coherent sheaf of ideals $\mathcal{I}$
there exists an $m > 0$ such that
$H^1(X, \mathcal{I} \otimes \mathcal{L}^{\otimes m})$ is zero.
This will finish the proof by
Cohomology of Schemes, Lemma \ref{coherent-lemma-vanshing-gives-ample}.
To see this we consider the maps
$$
\mathcal{I} \xrightarrow{s}
\mathcal{I} \otimes \mathcal{L}^{\otimes n} \xrightarrow{s}
\mathcal{I} \otimes \mathcal{L}^{\otimes 2n} \xrightarrow{s} \ldots
$$
Since $\mathcal{I}$ is torsion free, these maps are injective and
isomorphisms over $X_s$, hence the cokernels have vanishing $H^1$
(by Cohomology of Schemes, Lemma
\ref{coherent-lemma-coherent-support-dimension-0} for example).
We conclude that the maps of vector spaces
$$
H^1(X, \mathcal{I}) \to
H^1(X, \mathcal{I} \otimes \mathcal{L}^{\otimes n}) \to
H^1(X, \mathcal{I} \otimes \mathcal{L}^{\otimes 2n}) \to \ldots
$$
are surjective. On the other hand, the dimension of $H^1(X, \mathcal{I})$
is finite, and every element maps to zero eventually by
Cohomology of Schemes, Lemma
\ref{coherent-lemma-section-affine-open-kills-classes}.
Thus for some $e > 0$ we see that
$H^1(X, \mathcal{I} \otimes \mathcal{L}^{\otimes en})$ is zero.
This finishes the proof.
\end{proof}

\end{document}
